\documentclass[12pt]{article}
\usepackage{kyradefaults}
\addbibresource{references.bib}

\title{\textbf{All Aboard? Causal Evidence on the Labor-Market Effects of New Rail Stations}\\ \medskip
\large Interim Proposal for Topics in Social Insurance}
\author{Kyra Sadovi%
  \thanks{E-Mail: \texttt{ksadovi@uchicago.edu}}}
\affil{Harris School of Public Policy, University of Chicago}
\date{}

\begin{document}
\maketitle
\thispagestyle{empty}
 \begin{abstract}
     In this proposal, I set out the framing for a paper which estimates the causal effect of new rail station access on nearby residents' labor-market outcomes, focusing on worker flows and income. New stations may improve outcomes through three mechanisms: more reliable commuting to existing jobs, expanded job-search radius across the metro area, and broader referral networks. Identifying these effects is complicated by non-random station placement, residential sorting, and anticipatory adjustment. I address these challenges by exploiting quasi-random variation in station opening timing generated by construction delays, embedding the design in a stacked difference-in-differences framework that yields robust event-study estimates under staggered adoption. A novel dataset of U.S. rail projects since 2000 links planned opening dates — drawn from Draft Environmental Impact Statements — to realized opening dates and station geographies. Labor-market outcomes come from Census LODES, and transit exposure is measured via walking-time isochrones computed using the \texttt{r5r} routing package.
 \end{abstract}
\newpage 
\section{Research Question}
I want to measure the impact of increased access to public transit on surrounding residents' labor market outcomes. I propose using U.S. light rail project construction delays as a plausibly exogenous source of variation for a staggered difference-in-differences approach, thereby isolating the effect of gaining physical access to a transit station. 

It is important to emphasize that I am interested in the network effects of gaining access to a transit system on residents' outcomes rather than the overall, multifaceted impact of transit. The reason I make this distinction is because my delay setup intentionally differences out all anticipation effects. As \textcite{mcmillenReactionHousePrices2004} find in their paper on the opening of Chicago's Orange Line, non-labor economic indicators such as house prices can respond to the announcement of a planned station up to six years in advance of the station actually opening. While this is an interesting research question in and of itself, it is not the subject of this proposal or the subsequent paper. 

Decreased commuting costs impact labor market outcomes through three distinct (though related) mechanisms: more reliable means of commuting make it more likely that a worker retains their current employment \parencite{sanchezConnectionPublicTransit1999}; an expanded plausible job-search radius within a given metropolitan area \parencite{perezUrbanTransitInfrastructure2022, donaldsonRailroadsAmericanEconomic2016}; and broader referral networks \parencite{patacchini_referral_fx, kimSpatialInteractions2023}. My model setup is specifically designed to isolate transit's effects on residents via these mechanisms --- in short, to measure the effect of gaining access to a system rather than the effect of being chosen for a transit station. Comparing tracts with stations to those which are soon-to-receive stations allows me to estimate this effect, essentially the average treatment effect on the treated (with treatment here being receipt of a station). 

While I intend my reduced form approach to focus on labor market outcomes (specifically, my first stage is worker flows), a potential welfare analysis could extend the \textcite{borusyakNonrandomExposureExogenous2023} work on non-random exposure shocks, particularly using my delay setup. Not only would my extension provide a similar parameter estimate in the U.S. setting, but it would also expand on preexisting econometric approaches by combining the \citeauthor{borusyakNonrandomExposureExogenous2023} recentered instrument and a stacked difference-in-difference estimator (created by stacking by calendar year due to the project time indicator, explained in section \ref{sec:design}).

While this is a question that has been the subject of much of the transit economics literature (see \textcite{tsivanidisEvaluatingImpactUrban2026}, \textcite{tyndall_airports}, and \textcite{Gupta} as just a few examples), my setting and specific posture offers the potential to isolate transit impacts on agents as residents rather than as workers --- effectively, measuring the other node of the commute to work relative to \citeauthor{tsivanidisEvaluatingImpactUrban2026}. But even relative to \citeauthor{tyndall_airports}'s work using airport lines as an albeit ingenious instrument to measure labor market outcomes, my proposal offers the potential to expand on his case-study setup and make the results generalizable to transit projects across the country. 

\section{Institutional Setting}
If I were given carte blanche to measure the labor market or welfare impacts of public transit, I would have to completely randomize resident locations, workplace locations, and resident demographics as well as station location characteristics in order to isolate transit's impact from the many sources of endogeneity characteristic of analyses like this. As I am not afforded that convenience, I contest that my delay setup does a good job of isolating these effects. 

As I'll discuss more in section \ref{sec:data}, I have collected data on construction and environmental impact records of new transit projects in urban American cities in the last twenty years. The reason I've focused on the United States is because of the country's particular attitude toward public transportation and increasing trends of urbanization and congestion \parencite{roseTrafficCongestionHits2025, fhwa2026ucr}. Whether or not Americans are reticent to use public transit or are simply not afforded the opportunity to do so is the underlying debate among public policymakers deciding whether to invest more funds in public transit projects. Rejecting the null that Americans do not change their behavior when given access to a new transit station would support the claim that increased funding for public transit would increase residents' welfare and employment outcomes. While the same question could be asked in other settings (and has), those results are often considered not generalizable to the American worker. This paper could answer that concern directly. 

\section{Data \& Feasibility}\label{sec:data}
There are two sources of data I need to complete this analysis: information on public transit construction projects, and time-series data on labor-market outcomes at least at the Census tract level (the more granular, the better). 

The transit data that needs to be collected is all new station locations in the last 20 years; their opened dates; their announcement dates; and their projected opened dates. I have collected data on 230 stations so far, 94 of which have been opened. I started this collection based on the dataset compiled by \textcite{levy_2025_9wnjp-kez15}, but as I started filling in the delay information system by system I noticed a number of overlooked projects. I've since been updating this as I go, but currently am only running estimates with the \citeauthor{levy_2025_9wnjp-kez15} stations. Further, allowing for some time to pass for more recent stations to come online will benefit my number of observations greatly. I discuss this below in my MDE discussion. 

In an ideal world, it would be most useful to have longitudinal panel data on employment outcomes at the resident level over time. I would need information on where the person lives and where they work --- the type of data which is available in the confidential LEHD Origin-Destination Employment Statistics (LODES) data \parencite{census_lodes}. Being constrained to public data is not a fundamental challenge, however: the public LODES data is available at fairly granular geographic levels, and given the amount of variation available from the transit data I have it shouldn't be an issue for at least initial treatment effect estimates. 

However, there's a larger problem with LODES: data reporting discrepancies. As \textcite{abrahamExploringDifferencesEmployment2013} demonstrate, employers often misreport their employees' places of work. LODES and other Census data repositories tend to source their employment data from employers themselves rather than workers; as a result, employer-side data reporting issues arise rather than worker discrepancies. Employers, consequently, will either mistakenly or intentionally (for reporting burden reasons, most frequently) report that their employees' physical place of work is the headquarters site rather than their actual physical office. This poses fundamental data credibility issues for the LODES-sourced worker flow data. 

As a result, my ideal worker information data would be sourced from \textcite{infogroup_residential}, which is available via the Becker Friedman Institute. These data solve the granularity issue and would enable me to get microeconomic estimates of treatment effects based on a number of rich demographic data rather than tract aggregates. While it does not solve the issue of employer location misreporting, the fact that it can be matched with external data gives rise to potentially more credible data in the future. 

The MDE from the static TWFE is approximately 0.082 log points (8.6\%) for log inflows, derived from the cluster-robust standard error on open (clustered at the tract level) multiplied by 2.8 (the standard factor for $\alpha = 0.05, \text{ power} = 0.8$). This is a plausible effect size relative to the literature: \textcite{Gupta} find property value effects of roughly 5–10\% near new subway stations, and labor market responses could reasonably fall in a similar range. The sample incorporates all projects cataloged by  \textcite{levy_2025_9wnjp-kez15}, supplemented by additional projects identified independently that do not yet appear in that database. As a result, the station sample is being updated in real time and is likely to grow — both through the addition of newly identified historical projects and, importantly, as stations already logged in the dataset come online in the coming years, converting currently untreated observations into post-opening variation. One caveat is that the event study in event time $k$ uses a narrower control group (pre-opening observations from eventually-opened stations only, dropping never-opened stations), which may modestly reduce power relative to the static TWFE figure reported here. Overall, the design appears adequately powered to detect moderate treatment effects at the current sample size, with the expectation that power will improve as the dataset expands.

\section{Research Design}\label{sec:design}
The parameter of interest is the average treatment effect on the treated (ATT) of a new transit station opening on surrounding labor markets --- specifically, the causal effect on log worker inflows and outflows to census tracts that actually receive a station. I also seek to recover the dynamic version of this parameter: how the effect evolves over time relative to actual station opening, which requires identifying a separate coefficient for each period $k$ in event time.

Identification relies on the conditional exogeneity of construction delays. Let $$j = \text{census year} - \text{initial expected open year}$$ denote project time --- the year relative to when a station was originally projected to open, as recorded in the Draft Environmental Impact Statement. The key assumption is that, conditional on project time $j$ and census tract fixed effects, the timing of actual station opening is as-good-as-random:
$$\mathbb{E}[Y_{it}(0) \mid \text{open}_{it},\, j_{it},\, \text{tract}_i] = \mathbb{E}[Y_{it}(0) \mid j_{it},\, \text{tract}_i].$$
This is motivated by the institutional observation that construction delays are driven primarily by engineering complications, contracting failures, and subsurface geology rather than by anticipatory changes in neighborhood economic conditions. Critically, initial projected opening dates are fixed at the DEIS or FEIS stage --- years before construction begins --- meaning they predate any local economic response to the project.

Among the identifying assumptions, I view conditional parallel trends as the most credible, since project-time fixed effects absorb the systematic relationship between neighborhood characteristics and project schedules. The conditional exogeneity of delays is somewhat less certain, as political economy considerations could cause stations in economically distressed areas to face longer delays, though tract fixed effects mitigate this concern. The least credible assumption is SUTVA: tracts near multiple stations, or near both an opened and a delayed station, could experience spillovers that contaminate the control group.

The institutional features most important for identification are (1) that initial projected opening dates are fixed in federal planning documents well before construction begins, making them plausibly orthogonal to subsequent local trends, and (2) that the proximate causes of delay --- underground obstructions, cost overruns, contractor disputes --- are unrelated to the economic trajectory of adjacent neighborhoods. I provide several pieces of evidence in support of these features. Figure \ref{fig:balance_covariates} reports covariate balance across groups defined by delay length: point estimates are small relative to their confidence intervals and centered near zero across all six pre-treatment tract characteristics, consistent with delays being orthogonal to baseline neighborhood conditions. Figure \ref{fig:delay_income} plots construction delay against tract-level median household income and reveals no discernible relationship, further supporting the claim that longer delays are not concentrated in systematically richer or poorer neighborhoods. Figure \ref{fig:balance_regression} presents coefficients from a regression of delay length on pre-treatment tract characteristics; none of the covariates are precisely estimated and all confidence intervals include zero. Figures \ref{fig:delay_state} and \ref{fig:delay_system} additionally document that delay variation is distributed across states and transit systems rather than being driven by a single agency or geography, reducing concern that unobserved system-level confounders are driving the results.

The headline specification is a static two-way fixed effects regression,
$$\log Y_{it} = \beta\, \texttt{open}_{it} + \alpha_i + \delta_j + \varepsilon_{it},$$
where $\alpha_i$ are census tract fixed effects, $\delta_j$ are project-time fixed effects, and $\texttt{open}_{it} = \mathbf{1}[\text{census\_year} \geq \text{year}(\text{open\_date}_i)]$. For dynamic effects, I estimate an event study replacing \texttt{open} with a full set of dummies for $k = \text{census\_year} - \text{year}(\text{open\_date})$, normalizing to $k = -1$. Standard errors are clustered at the census tract level throughout, with two-way clustering at the tract and transit-system level as a robustness check.

Preliminary results are encouraging. Figure \ref{fig:es_inflows} presents the event study for log inflows. Pre-opening coefficients are flat and statistically indistinguishable from zero in the years immediately preceding opening, consistent with the parallel trends assumption. Following station opening, inflows rise steadily and persistently, reaching approximately one log point above baseline by ten years post-opening. The growing magnitude of the effect over time suggests that transit access reshapes residential sorting gradually, as workers update location decisions in response to improved commute options, rather than producing a one-time discrete jump at opening.
\newpage
\printbibliography
\newpage
\section*{Appendix}\label{sec:appendix}
\subsection*{Figures}\label{subsec:figs}
\begin{figure}[H]
    \centering
    \includegraphics[width=0.8\linewidth]{3_output/2_figures/3_delay_robustness/cross_vs_within_system_coef_plot.png}
    \caption{}\label{fig:balance_covariates}
\end{figure}
\begin{figure}[H]
    \centering
    \includegraphics[width=0.8\linewidth]{3_output/2_figures/3_delay_robustness/delay_vs_income.png}
    \caption{}\label{fig:delay_income}
\end{figure}
\begin{figure}[H]
    \centering
    \includegraphics[width=0.8\linewidth]{3_output/2_figures/3_delay_robustness/within_system_coef_plot.png}
    \caption{}
    \label{fig:balance_regression}
\end{figure}
\begin{figure}[H]
    \centering
    \includegraphics[width=0.8\linewidth]{3_output/2_figures/3_delay_robustness/mean_delay_by_state.png}
    \caption{}
    \label{fig:delay_state}
\end{figure}
\begin{figure}[H]
    \centering
    \includegraphics[width=0.8\linewidth]{3_output/2_figures/3_delay_robustness/mean_delay_by_system.png}
    \caption{}
    \label{fig:delay_system}
\end{figure}
\end{document}